\hypertarget{appendix-fecc1847f5a3}{}
\section*{Limits and continuity}
\chapter{Limits and continuity}\label{appendix-fecc1847f5a3-limits-and-continuity}

\textbf{Curriculum depth:} Developed · \textbf{First seen:} Phase 9 · \textbf{Latest development:} Phase 9

\section{Definition and mechanism}\label{appendix-fecc1847f5a3-definition-and-mechanism}

A limit describes the value a function approaches as its input approaches a point, even when the function is undefined or assigned a different value there. Continuity at a point adds the requirement that the function value equal that limit. Epsilon-delta language makes this precise: every requested output tolerance can be met by choosing a sufficiently small input neighborhood. Continuity rules out finite jumps but does not guarantee a derivative; the absolute-value function is continuous at zero yet has a corner there. Limits expose removable singularities such as a cancelled factor, while one-sided limits reveal boundaries and discontinuities. In numerical software, local smoothness matters because derivative-based approximations rely on nearby inputs producing controlled output changes. Domain restrictions, branch conditions, and singular denominators must be checked before applying derivative formulas.

\section{Prerequisites}\label{appendix-fecc1847f5a3-prerequisites}

\begin{itemize}
\tightlist
\item
  No earlier concept node is required in this planned path.
\end{itemize}

\section{Used by}\label{appendix-fecc1847f5a3-used-by}

\begin{itemize}
\tightlist
\item
  \hyperref[appendix-d9f70bdfbc92]{Differentiation rules and chain rule}
\end{itemize}

\section{Full teaching explanation}\label{appendix-fecc1847f5a3-full-teaching-explanation}

\begin{itemize}
\tightlist
\item
  \href{Continuous\%20Notes/Computing\%20Foundations\%20-\%20From\%20Bits\%20to\%20C++\%20and\%20Systems\%20APIs\#Functions,\%20limits,\%20continuity,\%20and\%20scalar\%20change}{Phase 9 teaching section}
\item
  \href{../01-Foundations-I/main.pdf\#nameddest=appendix-adf48f5dc5c1}{Visual Recall/09 - Calculus and Optimization - Visual Recall.md (01-Foundations-I, appendix-adf48f5dc5c1)} · visual recall
\end{itemize}

\section{Source coverage}\label{appendix-fecc1847f5a3-source-coverage}

Supplied lesson parts: 803 (T04 · M049), 804 (T04 · M049).
